\section*{The Propers: Themes and Movement}

This solemnity moves from God's self-disclosure to the Church's participation: mercy is proclaimed, the Son is given, the Spirit forms communion and filial speech, and worshippers return the gift in confession and offering. The unity is liturgical and doctrinal.

\begin{center}
\begin{tikzpicture}[x=3.35cm]
\node[guide node,text width=2.55cm] (mercy) at (0,0) {\textbf{Mercy disclosed}\\Name, covenant, pardon};
\node[guide node,text width=2.55cm] (gift) at (1,0) {\textbf{Gift enacted}\\Son sent for salvation};
\node[guide emphasis,text width=2.55cm] (communion) at (2,0) {\textbf{Communion given}\\Grace, love, Spirit};
\node[guide node,text width=2.55cm] (return) at (3,0) {\textbf{Gift returned}\\Offering, adoption, praise};
\foreach \x/\y in {mercy/gift,gift/communion,communion/return}
  \draw[guide arrow] (\x) -- (\y);
\end{tikzpicture}
\end{center}

\subsection*{Mercy Reveals the One God\enspace\elementref{Entrance, Collect, First Reading, Canticle}}

The non-scriptural Entrance Antiphon blesses Father, Son, and Spirit for merciful love before the Collect names the missions of Word and sanctifying Spirit. Exodus supplies the covenantal grammar beneath that praise: God freely proclaims the divine Name and qualities, while Moses responds with worship and a plea for pardoning presence. Daniel then makes the response doxological from within persecution. The movement is revelation, adoration, petition, and blessing---not a deduction of three Persons from repeated Old Testament phrases.

\begin{studybox}{Revelation, unity, and distinction}
The works by which God reveals and gives himself truly disclose the Trinity, yet the Old Testament does not thereby become an explicit Nicene treatise (Catechism 236--237). One divine nature and one common operation exclude three competing gods; real personal distinction excludes three temporary masks (Catechism 253--258).
\end{studybox}

\subsection*{The Gift of the Son Discloses Salvation and Judgment\enspace\elementref{Acclamation, Gospel}}

The adapted acclamation names the eternal Trinity before John proclaims the Father's gift of the Son. John holds purpose and crisis together: the Son is sent to save, not to seek condemnation, but the light's coming exposes the response of belief or refusal. Augustine's direct exposition keeps the medicinal image central. Salvation is offered through the crucified and exalted Son; judgment is not an arbitrary divine delight in ruin.

\subsection*{Apostolic Blessing Becomes Ecclesial Form\enspace\elementref{Second Reading}}

Paul's blessing does not float above the preceding imperatives. Grace, love, and communion crown a call to repair, encouragement, agreement, and peace. Chrysostom sees the three named sources without dividing divine gifts into separate economies. The solemnity therefore places dogmatic confession inside a concrete ecclesial task: a people blessed in Trinitarian form is to become reconciled and peaceable.

\begin{studybox}{Apostolic blessing and later dogma}
The apostolic blessing is authentic triadic revelation, but the technical vocabulary of Nicaea and Constantinople comes after Paul. Later dogma protects the apostolic witness from tritheist, modalist, and subordinationist misreadings.
\end{studybox}

\subsection*{Confession Passes into Offering and Adoption\enspace\elementref{Prayer over Offerings, Preface, Communion, Prayer after Communion}}

The Prayer over the Offerings asks that invocation of the divine Name sanctify both the gifts and the worshippers' enduring self-offering. The Preface supplies exact doctrinal grammar: distinct Persons, one substance, equal majesty. At Communion, Galatians makes confession lived relationship; the Spirit of the Son voices adopted address to the Father. The final prayer asks that sacramental reception and confessed faith profit body and soul. Doctrine, sacrament, and life remain distinguishable but inseparable.
